\section{Recomendaciones}

Con base en los resultados obtenidos durante la práctica, el Grupo 6 propone una serie de recomendaciones orientadas a fortalecer la seguridad de redes y sistemas informáticos:

\begin{itemize}
    \item Implementar políticas de segmentación de red más estrictas mediante VLANs y listas de control de acceso (ACL), con el fin de limitar la comunicación únicamente a los servicios necesarios.
    
    \item Restringir y monitorear los puertos abiertos en los servidores, asegurando que únicamente aquellos indispensables para el funcionamiento del sistema permanezcan activos.
    
    \item Incorporar mecanismos de protección contra ataques de denegación de servicio, tales como firewalls, sistemas de detección y prevención de intrusiones (IDS/IPS) y limitación de tráfico.
    
    \item Fortalecer la seguridad de las redes inalámbricas mediante el uso de cifrado robusto (WPA3), contraseñas complejas y la desactivación de funciones vulnerables.
    
    \item Evitar la ejecución de archivos de procedencia desconocida en sistemas Windows, complementando esta medida con soluciones antivirus y antimalware actualizadas.
    
    \item Mantener los sistemas operativos, dispositivos de red y aplicaciones constantemente actualizados para reducir la exposición a vulnerabilidades conocidas.
    
    \item Capacitar a los usuarios en buenas prácticas de ciberseguridad, enfatizando los riesgos asociados a la ingeniería social y a la manipulación de archivos maliciosos.
\end{itemize}

La aplicación de estas recomendaciones contribuirá a minimizar los riesgos de ataques informáticos y a mejorar la postura de seguridad de las organizaciones, permitiendo una gestión más segura y eficiente de sus infraestructuras tecnológicas.
